We use Andrew's monotone chain algorithm~\citep{andrew1979another} to compute
the convex hull of a district polygon in $O(n \log n)$ time.
(Graham scan~\citep{graham1972efficient} achieves the same asymptotic complexity
but requires a polar angle sort; Andrew's algorithm uses a lexicographic sort
that is numerically more stable.)

\subsection{Algorithm}

\begin{algorithm}
\caption{Andrew's Monotone Chain Convex Hull}
\label{alg:monotone}
\begin{algorithmic}[1]
\Function{ConvexHull}{$P$: sorted point set}
  \State Sort $P$ lexicographically by $(x, y)$
  \State Build lower hull $L$:
  \For{$p \in P$}
    \While{$|L| \geq 2$ and cross$(L[-2], L[-1], p) \leq 0$}
      \State Remove last point from $L$
    \EndWhile
    \State Append $p$ to $L$
  \EndFor
  \State Build upper hull $U$ similarly (iterate $P$ in reverse)
  \State \Return $L \cup U$ (removing duplicate endpoints)
\EndFunction
\end{algorithmic}
\end{algorithm}

The \texttt{cross} function computes the cross product:
\[
  \text{cross}(O, A, B) = (A_x - O_x)(B_y - O_y) - (A_y - O_y)(B_x - O_x).
\]
A non-positive cross product indicates a right turn (or collinear point),
which is removed to maintain convexity.

\subsection{Convex Hull Area}

After computing the convex hull vertices, the area is computed via the shoelace formula:
\[
  A(\text{conv}(D)) = \frac{1}{2} \left| \sum_{i=0}^{m-1} (x_i y_{i+1} - x_{i+1} y_i) \right|,
\]
where $(x_i, y_i)$ are the vertices of the convex hull polygon.

\subsection{Test Invariants}

L0 invariants:
\begin{itemize}
  \item CH of any convex polygon is exactly $1.0$.
  \item CH $\in (0, 1]$ for any polygon.
  \item For a cross-shaped synthetic district with arm length 10 and width 2:
    CH $\approx 0.80$ (verifiable analytically).
  \item Convex hull area $\geq$ district area for all test geometries.
\end{itemize}
